\chapter*{Introduzione}

I Large Language Model hanno reso più accessibile l'interazione con grandi
quantità di testo: possono sintetizzare documenti, formulare risposte in
linguaggio naturale e assistere attività di scrittura o analisi. In ambito
aziendale, però, queste capacità non bastano. Una risposta utile deve riferirsi
a dati interni aggiornati, deve essere verificabile e deve rispettare vincoli di
riservatezza, controllo dell'infrastruttura e manutenibilità. Un modello
pre-addestrato, usato da solo, non conosce la documentazione privata di
un'organizzazione e non possiede un meccanismo nativo per controllarla. Il
problema emerge soprattutto quando il sistema deve rispondere su contenuti
specifici dell'azienda, non su conoscenza generale.

Il paradigma \emph{Retrieval-Augmented Generation} (RAG) affronta questo
problema combinando due componenti: un modulo di recupero, che seleziona
documenti pertinenti da una base di conoscenza esterna, e un modello
generativo, che produce la risposta a partire dalla domanda e dal contesto
recuperato~\cite{lewis2020rag}. In questo modo la conoscenza usata dal sistema
non risiede soltanto nei parametri del modello, ma anche in un indice
aggiornabile e ispezionabile. Per scenari enterprise, tale impostazione è
interessante perché consente di integrare fonti interne, mantenere traccia dei
documenti usati e separare il problema della generazione dal problema della
gestione della conoscenza.

Questa tesi nasce all'interno di Intré S.r.l., una software factory italiana
che adotta pratiche strutturate di apprendimento organizzativo. Tra queste,
le Gilde rappresentano cicli periodici di apprendimento di gruppo: le persone
propongono temi, formano gruppi di lavoro, producono risultati e condividono
l'esperienza con il resto dell'azienda. Nel tempo, la piattaforma interna
i3Guild ha raccolto proposte, descrizioni, partecipanti, risultati e materiali
collegati. Tale corpus costituisce una memoria organizzativa di valore, ma non
sempre facile da interrogare: le informazioni sono distribuite tra database,
metadati e contenuti testuali, e il recupero manuale richiede conoscenza
pregressa del dominio.

L'obiettivo principale della tesi è progettare, implementare e valutare un
sistema RAG integrato in i3Guild, capace di rendere interrogabile il corpus
delle Gilde tramite linguaggio naturale. Il sistema deve rispettare tre vincoli
di fondo. Il primo è la riservatezza: i dati aziendali devono restare entro un
perimetro controllato, in coerenza con le esigenze di protezione del dato e con
il quadro normativo europeo~\cite{gdpr2016}. Il secondo è la sostenibilità
infrastrutturale: il prototipo deve poter funzionare con risorse realistiche,
senza assumere la disponibilità di GPU server dedicate. Il terzo è la
manutenibilità: la soluzione deve integrarsi con l'applicazione esistente e
conservarne la separazione tra sorgente dati, indicizzazione, retrieval e
generazione.

Il lavoro si concentra sull'integrazione di componenti
esistenti in un prototipo enterprise on-premise: PostgreSQL come sorgente
autoritativa, una pipeline di ingestione per normalizzare i dati e produrre
embedding, Qdrant come database vettoriale, Haystack come framework di
orchestrazione, FastAPI come microservizio RAG, Ollama come server di inferenza
locale e Next.js come livello applicativo di integrazione. La parte
implementativa include inoltre query rewriting, filtri strutturati,
deduplicazione, selezione dinamica dei documenti, streaming della risposta e
generazione controllata di bozze di nuove Gilde.

Durante lo sviluppo è stato esplorato anche un workflow agentico per guidare
l'utente nella proposta di nuove Gilde. Tale workflow è stato implementato e
testato, ma non adottato come nucleo della soluzione finale. La ragione è
pratica: a parità di hardware locale, l'approccio agentico introduce più
passaggi, maggiore costo computazionale e maggiore complessità operativa. Nel
perimetro del prototipo, un RAG più classico e deterministico è risultato più
stabile e più semplice da mantenere. La componente agentica rimane quindi
documentata come sperimentazione e possibile sviluppo futuro.

La seconda parte del lavoro affronta un problema emerso dal prototipo: la
sostenibilità dell'inferenza locale. Un sistema RAG aziendale può essere
correttamente progettato, ma la sua utilità pratica dipende anche dal modello
generativo usato nella fase finale. Il modello deve essere
caricabile, deve rientrare nella memoria disponibile, deve rispondere con una
latenza accettabile e deve mantenere un livello qualitativo compatibile con il
caso d'uso. Per questo la tesi include un capitolo sperimentale dedicato
all'inferenza locale quantizzata su Apple Silicon, usando MLX e modelli
open-weight tra 3B e 14B parametri.

L'esperimento misura il compromesso introdotto dalla quantizzazione MLX su un
MacBook Pro con chip M1 Pro e 16 GB di memoria unificata. Le varianti FP16,
Q8, Q6 e Q4 vengono confrontate rispetto a dimensione su disco, memoria di
picco, tempo di caricamento, time-to-first-token, throughput di decodifica e
qualità su subset ridotti MMLU e GSM8K. I risultati mostrano che, per i
modelli piccoli, Q4 riduce la dimensione di circa il 72\%, riduce il picco di
memoria di circa 67--70\% e aumenta il throughput di decodifica di circa
2.6--2.8 volte rispetto a FP16, senza mostrare un degrado macroscopico nei
controlli automatici ridotti. Per i modelli medi e grandi, la quantizzazione è
soprattutto una leva di fattibilità: rende eseguibili configurazioni che
sarebbero meno realistiche nel vincolo dei 16 GB, ma richiede una valutazione
attenta di latenza, formato dell'output e uso concorrente con gli altri servizi
del sistema.

La tesi è organizzata come segue. Il Capitolo~\ref{chap:background} introduce
i fondamenti tecnologici: Transformer, LLM, inferenza locale, RAG, embedding,
database vettoriali, Agentic RAG e Haystack. Il
Capitolo~\ref{chap:requirements_design} analizza il contesto aziendale di
Intré, il modello delle Gilde, il caso d'uso, i requisiti e le scelte
architetturali. Il Capitolo~\ref{chap:implementazione} descrive
l'implementazione del sistema: pipeline di ingestione, microservizio RAG,
retrieval, generazione locale, workflow agentico sperimentale e integrazione
Next.js. Il Capitolo~\ref{chap:ottimizzazione} definisce il protocollo
sperimentale per l'inferenza locale quantizzata su Apple Silicon. Il
Capitolo~\ref{chap:evaluation} presenta e discute i risultati, collegandoli
alle implicazioni operative per un sistema RAG locale.
